\documentclass[12pt,a4paper]{ctexart}

% Packages
\usepackage{amsmath,amssymb}
\usepackage{graphicx}
\usepackage{booktabs}
\usepackage{hyperref}
\usepackage{geometry}
\usepackage{caption}
\usepackage{float}

\geometry{left=2.5cm,right=2.5cm,top=2.5cm,bottom=2.5cm}

\graphicspath{{figures/}}

% Metadata
\title{基于多源数据融合的智能天气分析系统设计与实现}
\author{徐毅成}
\date{\today}

\begin{document}

\maketitle

\begin{abstract}
随着气象信息服务从原来的单一预报展示发展到现在的综合分析、智能辅助和行业应用，
用户对天气平台的要求也由原来的只关心温度、降水等基本信息，发展为需要系统提供更加细致的预报对比、
更加专业的大气诊断、更加接近实际场景的决策支持等。本文根据以上需求设计出一套基于Flask的智能气象信息服务系统。
系统主要面向杭州地区，以多模式预报分析为主线，集成了数值预报展示、探空分析、历史气候分析、
机器学习误差校正、农业气象服务等功能模块，形成了一种综合性、实用性较强的Web应用平台。

本文用 Flask 构建了后端服务框架，并用 Pandas、NumPy 等工具对气象数据进行处理。
系统在预报分析模块里把欧洲中期天气预报中心全球模式（ECMWF）、美国全球预报系统（GFS）、
德国气象局全球模式（ICON）等模式预报结果融合起来，给出 72 小时精细化预报及多模式对比显示；
在探空分析模块里完成了怀俄明大学探空资料抓取、层结解析、稳定度参数算出并做图；
历史气候模块可以做年度气候统计、趋势分析，基于 EW 指数的代表性极端事件挑选；
机器学习误差校正模块把 ECMWF/Open-Meteo 的预报样本做了站点口径统一、时次统一、模型重训；
农业气象模块创建了作物知识库、积温计算、农事建议等模块，表现出系统的行业拓展能力。

本文主要论述了系统的需要分析、总体架构设计以及重要模块的实现方法，
并通过页面功能、接口流程、模型评价结果等对系统进行测试分析。
结果表明，该系统可以较好的完成综合气象信息展示和分析的任务，
在预报分析的准确性、完整性以及应用的扩展性上都有一定的实践意义。

\textbf{关键词：}多源数据；智能气象服务；多模式预报；探空分析；历史气候；机器学习误差校正；农业应用
\end{abstract}

\tableofcontents
\newpage

% Include all chapters
\section{第 1 章 绪论}

\subsection{1.1 研究背景}

近些年来，由于数值天气预报技术、开放气象数据服务和Web应用开发技术的迅速发展，公众获取天气信息的方式也发生了明显的变化。早期的天气服务主要提供静态预报结果，即未来几天的气温、降水、风力等级等基本气象信息，而目前用户对气象产品的需要已经从原来的静态预报结果发展为现在的精细化、对比化、场景化的气象产品需求。一方面普通用户希望得到更加直观、更加连续的逐小时天气变化，从而为出行、生活提供帮助；另一方面，具有一定专业需求的用户，例如农业生产者、气象爱好者或者相关研究人员，更加关心不同模式之间存在的差异、探空层结结构、历史同类天气事件的统计特征以及预报误差的订正能力。

与此同时，开放气象数据源的丰富也给综合气象平台的创建打下了基础。以Open-Meteo为代表的开放预报接口使多模式预报结果的获取、整合变得更为方便；怀俄明大学提供的探空资料给大气稳定度分析、强对流诊断打下了基础；历史站点数据、再分析资料为气候趋势研究、极端事件识别提供长期样本。怎样把各种来源、时间跨度不同、业务逻辑各异的数据整合到一个平台上，用怎样的架构和交互方式向用户展示出来，成为了一个有实际操作价值的系统设计问题。

在此背景下，设计出一种可以满足普通用户需求，同时又具备一定专业分析能力的综合气象信息系统，有较高的应用价值和工程意义。该系统不但可以给杭州地区天气分析提供更加完整的工具链，还可以为之后扩展到更多的区域、更多的业务场景打下基础。

\subsection{1.2 研究意义}

本文的研究意义主要体现在以下几个方面。

在实际运用中，综合气象服务系统能克服单个预报展示平台所存在的缺陷。传统的天气页面只显示单个模式的结果或者很少的指标，不能反映出不同的模式之间存在的差异，也不能给探空、大气稳定度、历史同日极端性等更专业的内容提供支持。本文所实现的系统以多模式预报为主，把探空分析、历史统计、机器学习后处理整合在一个平台上，使用户可以在同一个界面里完成从短期预报浏览到深层次分析的各个步骤。

其次从工程实现角度出发，本文给出一个面向实际业务的中小型综合气象平台架构方案。系统不是简单的堆砌，而是以预报分析能力的建立为主线，把数据抓取、时间规整、站点配置统一、模型调用、分析计算和页面展示组织成比较清晰的模块结构。该种组织方式可以给同类型信息服务系统开发提供一定的借鉴。

再次从方法上来说，本文用历史同月同日的经验百分位构成EW指数来辅助识别出更具有代表性的一些极端天气事件，同时也对ECMWF/Open-Meteo预报数据进行了误差校正模型训练，将训练和运行时的站点口径、时次规则统一起来。这些工作没有追求复杂的算法革新，但是从系统实用性、结果可信度来说是有价值的。

最后在应用延伸方面，系统会把农业气象服务当作重要的扩展模块，尝试将气象分析能力同作物生长需求结合起来，创建起作物积温估算、风险提醒以及农事建议等功效，表现出综合气象平台向行业应用落地的可能。

\subsection{1.3 国内外研究与应用现状}

就国外的发展情况而言，气象服务平台已经普遍具有比较完善的在线化、可视化能力。部分业务平台可以提供多模式预报浏览、探空图显示、雷达和卫星叠加展示、专业参数分析等服务，说明气象服务正由原来的文本、表格形式的输出方式向交互式的平台转变。随着机器学习在数值预报后处理领域应用的增多，用历史预报和实况样本做偏差订正、概率预报和分类识别，已经成为提高预报质量的一个重要途径。

在国内，公共天气服务平台以及业务化气象平台近些年来也正在快速的发展，在精细化预报、短临预警、灾害监测以及行业服务等方面取得了较多的成果。但是对毕业设计或者中小型综合系统来说，常常会存在两个问题，一是平台的功能很容易停留在单个查询或者单个模块展示上，缺少多源数据融合之后的整体设计；二是即使系统包含了多个模块，也会出现功能堆砌的情况，缺少清晰的业务主线以及统一的数据组织逻辑。

因此，如何以“预报分析”为核心，兼顾历史气候、探空诊断、误差校正和农业服务等能力，构建一套结构清晰、实现完整、具备一定专业特色的综合气象信息服务系统，仍具有较强的实践研究价值。

\subsection{1.4 本文研究内容}

围绕上述目标，本文主要完成以下工作：

1. 设计并实现一个基于 Flask 的综合气象信息服务系统，构建统一的页面入口、接口组织和模块协同机制。

2. 建立多模式预报分析模块，集成 ECMWF、GFS、ICON 等模式数据，提供 72 小时精细化预报和多模式对比分析能力。

3. 建立探空分析模块，实现探空资料抓取、层结解析、稳定度参数计算和多种图形展示。

4. 建立历史气候分析模块，完成年度气候统计、趋势分析以及基于 EW 指数的代表性极端事件筛选。

5. 建立机器学习误差校正模块，对 ECMWF/Open-Meteo 预报结果进行站点统一、时次统一和模型评估。

6. 建立农业气象服务模块，将预报与历史分析能力延伸到作物生长和农事建议场景中。

需要说明的是，本课题开题阶段主要围绕综合天气展示、预报分析和历史气象等基础能力展开，农业气象服务并未作为独立模块被明确纳入正式开题边界。随着系统开发逐步推进，本文在完成核心预报分析能力的基础上，进一步扩展形成了农业气象服务链路。因而，论文正文以当前实际完成的系统为主要依据，在总体研究方向保持一致的前提下，对模块重点和结构安排作了适应性调整，即以预报分析为主线，以农业应用作为拓展内容展开论述。

\subsection{1.5 论文组织结构}

本文共分为七章。第一章为绪论，介绍研究背景、研究意义、国内外现状以及本文的主要工作。第二章介绍系统实现涉及的相关技术与理论基础。第三章对系统需求进行分析。第四章给出系统总体设计。第五章详细阐述各关键功能模块的设计与实现。第六章对系统进行测试与结果分析。第七章总结全文工作，并对后续改进方向进行展望。


\section{第 2 章 相关技术与理论基础}

\subsection{2.1 Flask 与前后端集成式 Web 架构}

本系统使用的是 Flask 作为后端的核心框架。相比于更重型的全栈框架，Flask的优势在于结构轻量、路由组织清晰、与模板页面和独立接口结合灵活，适合本项目这样“页面展示+数据接口+分析逻辑并存”的综合型应用。系统中的主页、多模式预报页、探空分析页、历史气候页和农业气象页均由 Flask 路由统一组织，而数据接口如 /api/advanced-forecast/multi-model、/upperair/data、/apply-ml-correction 等则承担前后端交互和异步数据更新功能。

该架构的业务意义就在于，它既保持了服务端模板渲染适合页面级整合展示的优点，又使关键模块可以独立地通过异步接口来更新数据，从而避免整页刷新造成的等待以及耦合问题。多模式预报、探空分析等都必须在页面加载之后再请求数据，这样的模板首屏+接口补充模式更适合气象信息平台的实时交互。

从系统组织角度来说，Flask 对于将路由层、服务层、分析层、模型层进行拆分是十分友好的。路由层主要是负责页面入口和接口调度，服务层主要是对外部数据进行获取和缓存组织，分析层主要是对历史气候、探空参数和农业指标进行计算，模型层主要是加载并调用机器学习误差校正模型。分层并不追求复杂的架构形式，但是可以很好地满足目前系统扩展的需求。

\subsection{2.2 气象数据来源与特点}

本系统使用的数据来源具有明显的多源异构特征，不同模块依赖的数据在时间尺度、空间尺度和更新机制上均存在差异，因此理解这些数据源的特点，是后续系统实现和结果解释的基础。

在数值预报方面，系统主要依赖 Open-Meteo 提供的开放接口获取 ECMWF、GFS、ICON 等模式结果。其优点在于调用便捷、格式统一、适合快速构建多模式对比能力。对本系统而言，Open-Meteo 不只是一个“拿天气数据”的接口，更是多模式预报分析和 ECMWF 后处理应用的基础输入源。

历史气象和农业积温用 ERA5-Land 再分析资料作为长期历史背景数据源。再分析资料的时间连续性好、历史资料丰富、字段统一，适合用作历史趋势分析、长期气候背景描述和积温这类长期积分型指标的估计。尽管再分析数据不是地面单站观测，但是在长期累积和区域背景分析方面有较好的稳定性以及可用性。

系统用怀俄明大学在线探空资料服务做高空观测数据的来源。该数据源可以给出标准探空时次的层结表以及指数资料，是做Skew-T分析、大气稳定度判断和对流潜势诊断的基础。与地面观测数据相比，探空资料可以更直接地反映出垂直方向上的温湿风结构，在系统中起到专业的诊断作用。

观测真值和机器学习误差校正方面，目前系统首先统一到杭州国家站/馒头山口径，如果更高的权威和稳定的本地小时观测获取有困难，那么阶段性使用Meteostat 58457作为当前最好的可用的统一真值源。虽然存在一定的工程折中，但是比训练、验证、运行时混用不同站点口径更合理，也更容易保证模型评价结果的可解释性。

\subsection{2.3 数据处理与分析支撑技术}

本系统在数据处理上主要使用了Pandas和NumPy这两个库。Pandas主要用在时间序列数据解析、缺测值清洗、按日或者按小时聚合、统计汇总、结构化输出上，NumPy更多地用来做数值计算、百分位统计、趋势计算、特征构造等。对气象系统来说，组合起来很适合处理时间字段多、缺测情况复杂、统计需求频繁的数据形态。

历史气候模块的年统计、月统计、同月同日样本提取、EW指数计算、季节事件筛选等都是用Pandas和NumPy完成的；农业模块中积温累计、适宜度评分、净有效降水等指标的计算也都基于这些基础处理能力之上。因此可以认为这两个工具是系统分析层最根本的数据处理基础。

另外系统在运行的时候还会加入简化的缓存机制来减少高频的数据请求造成的性能损失。多模式预报结果、模块状态以及部分提示信息都会被缓存起来，在短时间内不会对外部数据源进行重复请求。该种方式虽然不是复杂的分布式缓存方案，但是对目前规模的毕业设计系统来说，在保证时效性的同时也改善了响应体验。

\subsection{2.4 可视化与页面交互技术}

气象平台的结果表达高度依赖图形可视化，因此本系统在不同模块中采用了不同的图表技术组合。

在多模式预报和历史气候模块中，系统主要使用 Plotly.js 以及后端 Plotly / Plotly Express 生成交互式时间序列图、月度对比图、趋势图和分布图。Plotly 的优势在于交互能力较强，适合展示温度、降水、风速等连续变化过程，也方便在同一坐标系内对比多个模式或多个年份。

在探空分析模块中，系统主要使用 Matplotlib 与 MetPy。其中MetPy给出气象分析常用的绘图、热力学计算的支持，适合T-lnP图、Skew-T图、风矢量图等专业的图形绘制，Matplotlib是底层的绘图支持，保证探空图输出具有灵活性以及专业性。对于该模块来说，选择Matplotlib+MetPy比通用Web图表库更适合，因为探空图本身就是一种专业性的气象图形。

农业模块页面布局主要是用 Bootstrap，专题页环境曲线用 Chart.js 绘制。bootstrap适合于组织多卡片、多区域的信息页面，Chart.js适合绘制较轻量的专题曲线图。前端交互部分用JavaScript实现异步数据请求和局部刷新，农业页面中AI建议刷新、探空页面的图形切换等都属于此类。

该种“不同的模块使用不同的图形技术”是符合系统实际需要的。也就是说，本系统并没有因为统一技术栈而强行使用一种图表库，而是根据业务场景选择更加适合的表达方式。

\subsection{2.5 常用气象分析方法基础}

本系统虽然不以理论创新为目标，但仍然用到了若干具有明确气象业务背景的分析方法。

在多模式预报分析中，核心思想是通过同一时空位置下不同模式结果的并列展示，帮助用户判断天气过程强度、转折时间和趋势演变的一致性或分歧程度。这种方法本质上属于数值预报结果的比较分析，强调“同一事件、不同模式”的横向判断。

在探空分析中，系统使用 CAPE、CIN、K 指数、Lifted Index、整层可降水量和垂直风切变等指标来辅助判断大气不稳定条件和潜在风险。其中，Skew-T 或 T-lnP 图是分析温湿层结、抬升凝结和对流发展环境的经典图形工具，也是专业气象分析中的重要基础。

历史气候模块里，系统除了保留传统的原始极值排序之外，又加入了以历史同月同日的经验百分位为基础的EW指数，用来衡量某一日天气事件在长期历史背景下所具有的相对异常程度。对低温使用反向百分位、连续事件过程去重、四季分组展示，都是为了使极端事件识别更符合气候分析的业务逻辑。

农业模块中积温计算、适宜度评分属于比较典型的分析手段。积温用基点温度以上日平均有效热量的累计来表示作物发育进程，适宜度评分综合考虑温度、水分、极端天气的影响，对当前环境是否有利于作物生长做出定量的描述。虽然这些方法使用了简化的实现方式，但是已经可以满足农业管理提示的需要。

\subsection{2.6 机器学习误差校正基础}

机器学习误差校正在本系统中的定位，不是构建一个独立研究型算法平台，而是作为 ECMWF 预报链路的后处理增强层。当前系统采用 scikit-learn 中的随机森林模型完成温度、湿度、风速和降水的偏差订正，其中温度、湿度、风速使用回归模型，降水采用“先分类、后回归”的两阶段方案。

这一设计与业务问题本身是匹配的。温度、湿度和风速属于连续变量，更适合直接回归订正；降水则既包含“下不下”的分类问题，也包含“下多少”的量级问题，因此拆分建模更符合实际情况。特征方面，系统主要使用原始预报值、小时、月份、星期、昼夜属性、时效长度 lead\_hours、时效分桶以及周期性时间编码等变量，目的不是追求复杂模型，而是构建与运行时接口兼容、同时具备基本解释力的订正方案。

系统用issue\_time组织的滚动验证和最终留出测试来代替随机拆分所有目标时刻的方法。这样可以减少信息泄露，更接近真实业务中用过去起报样本预测未来起报样本的使用场景。本论文中这一部分的重要意义不单是模型本身，也是它体现出来的系统结果可信度以及工程闭环的提高。

\subsection{2.7 机器学习误差校正相关算法}

为了更清楚地说明本系统中误差校正模块的建模思路，本文对V1和V2两版方案中涉及到的主要算法进行简要说明。前者主要是用随机森林做回归和分类，后者也使用梯度提升树中的CatBoost做对比实验。

（1）随机森林(Random Forest)。随机森林属于一种基于决策树的集成学习方法，它主要依靠Bootstrap重采样来创建多个训练子集，然后单独训练多棵决策树，最后在预测时将各个树的结果加以整合。回归任务中模型一般输出各个树的结果平均值，分类任务用多数投票法得到最终类别。该方法抗过拟合能力强，能处理非线性关系，对异常值和局部噪声不敏感，适合当前样本量小、业务特征强的气象误差校正场景，所以本文V1模型使用了该算法。

（2）梯度提升树（CatBoost）。梯度提升树是通过不断训练新的树来拟合前一轮模型的残差，从而不断减小整体预测误差的一种方法。CatBoost是梯度提升树的高效工程实现，训练速度快、类别特征支持好、泛化能力强。在本文中，V2实验模型试图用CatBoost对ECMWF误差校正任务进行重训，以验证在统一站点口径、统一issue\_time规则、严格时间隔离评估条件下，boosting类模型能否得到比V1更好的结果。


\section{第 3 章 系统需求分析}

\subsection{3.1 总体需求概述}

本系统面向的并不是单一类型用户，而是同时覆盖普通天气查询用户、具有一定专业分析需求的用户以及农业应用场景下的业务型用户。因此，系统需求既包括基础的信息展示需求，也包括更深层的分析需求和一定程度的场景化决策支持需求。为了避免系统在功能组织上出现分散和割裂，本文将“预报分析能力”作为主线，在此基础上扩展历史统计、探空诊断、误差校正和农业服务等模块。

\subsection{3.2 用户需求分析}

普通用户最根本的需求就是得到迅速、清楚、可信、连续的天气信息。因此系统需要有易于理解的页面入口、当前天气概览、未来数小时变化趋势、未来数天变化趋势，在视觉上尽可能做到信息集中、切换方便、内容直观。

对于具备一定专业需求的用户而言，单一模式的天气结果往往不足以支持判断。他们更希望看到不同模式之间的对比关系、预报起报时次的合理解释、上空层结结构与稳定度参数，以及历史上同类型天气在多年样本中的相对异常程度。这类需求决定了系统不能只停留在“数据展示”，而要具备一定分析和解释能力。

对于农业场景用户而言，仅仅知道未来降水或温度并不够，还需要将天气信息与作物生长过程、积温进度和农业风险联系起来。因此，系统还应具备面向作物的积温估算、灾害风险提示和农事建议能力。

\subsection{3.3 功能需求分析}

根据目前系统的实现情况，本文把功能需求归纳为以下几方面。

第一，系统要具备多模式预报分析功能。该功能包含72小时精细化逐小时预报、多模式对比分析、预报趋势展示和ECMWF预报机器学习校正结果接入等，是整个平台的主要功能。

第二，系统要有探空分析功能。此功能可以对指定的时次进行探空数据抓取，得到层结表解析结果，绘制 Skew-T 和参数图，计算 CAPE、CIN、K 等稳定度参数。当抓取失败、时次不一致或者参数异常的时候，也应当具有一定的回退和异常提示功能。

第三，系统应该有历史气候分析的功能。该功能需要具备年度页、趋势页的展示能力，可以给出气温、降水等指标的长期统计结果，根据原始极值和EW指数分别代表事件两个方面来分析某一年里具有代表性的异常天气过程。

第四，系统要具有农业气象服务功能。该功能要依靠作物知识库，给典型作物提供积温进度、成熟日期估算、农业气象风险提醒以及农事建议。

第五，系统要有导航以及综合展示的功能。首页导航页要担负起平台总入口、主要模块分发、实时天气概览以及智能提醒汇总等任务，让用户在同一个页面上先建立起对于当前天气和系统能力结构的总体认识，然后进入相应的专业模块进行查看。

\subsection{3.4 非功能需求分析}

在非功能需求方面，系统至少需要满足以下要求。

响应性为第一。因为系统包含了很多外部数据源，如果没有缓存、并行请求或者失败回退等机制，那么页面加载就会出现明显迟滞。因此系统要依靠缓存、并发抓取以及接口分层来保证主要页面的可用性。

第二，稳定性。系统外部数据源有开放预报接口、探空资料源、历史观测数据源，不同数据源响应格式、可用时次、异常情况各不相同，所以系统要具有较强异常处理、容错能力。

第三，具有可扩展性。系统后续会增加更多的站点、更多的预报模式、更多的作物或者更复杂的机器学习模型，因此在架构设计的时候不能过于耦合，尽量使用分模块的方式。

第四，可以维护。由于系统包含页面模板、接口逻辑、数据抓取、分析计算、模型加载等多种代码，所以要保证模块的职责尽量清晰，便于以后的迭代。

\subsection{3.5 业务主线界定}

为了防止系统变成单纯的项目功能清单，本文在需求分析阶段就将预报分析能力的建立确定为业务主线。多模式预报模块直接给出未来天气分析能力，探空分析模块加强了专业的诊断能力，历史气候模块给目前的年份以及某些天气过程赋予了长时间的背景支撑，机器学习误差校正模块对预报品质展开后处理改良，农业气象模块在此基础上扩展到面向行业应用。这样一种组织方式既可以保证系统的综合性，又不会造成各个模块之间的相互隔离。


\section{第 4 章 系统总体设计}

\subsection{4.1 系统总体架构}

\begin{figure}[htbp]
\centering
\caption{图 4-1  智能气象信息服务系统总体架构图}
\label{fig:4-1}
\end{figure}

\begin{figure}[htbp]
\centering
\caption{图 4-2  系统核心数据流图}
\label{fig:4-2}
\end{figure}

\begin{figure}[htbp]
\centering
\caption{图 4-3  系统主要模块关系图}
\label{fig:4-3}
\end{figure}

\begin{figure}[htbp]
\centering
\caption{图 4-4  首页导航页界面示意图1}
\label{fig:4-4}
\end{figure}

\begin{figure}[htbp]
\centering
\caption{图 4-5  首页导航页界面示意图2}
\label{fig:4-5}
\end{figure}

本系统使用基于Flask的Web应用架构，整体可以分为路由层、服务层、分析层、模型层、展示层五个部分。路由层主要是处理页面入口以及接口调度，服务层主要是对外部数据进行抓取、整合以及运行时的配置，分析层主要是对历史气候、探空参数、农业指标等进行计算逻辑的处理，模型层主要是机器学习误差校正模型的加载和使用，展示层主要是由HTML模板、前端脚本、图表组件组成，实现数据可视化以及交互展示。

app.py 是系统主入口，主要用来组织首页、多模式预报、探空分析、历史气候、农业气象等各个页面的路由，并对相关的接口进行统一管理。以该入口文件为基准，系统把不同的业务逻辑拆分成独立的模块，多模式预报用advanced\_forecast\_service.py来处理，历史气候分析用history\_analyzer.py来处理，探空分析用sounding\_parser.py、sounding\_analyzer.py和sounding\_plotter.py协同处理，机器学习误差校正用ml\_correction.py和train\_bias\_correction.py支持，农业气象能力由作物知识库和相关计算模块提供。

\subsection{4.2 模块划分与职责}

系统是由五个主要业务模块构成的。

第一部分是多模式预报模块。该模块从Open-Meteo等处获取多种数值预报模式的结果，对逐小时预报进行组织、统一字段映射、预报时次估计以及缓存处理，最后向前端提供可以直接展示的数据结构。

第二，探空分析模块。从怀俄明大学探空接口抓取指定站点、指定时次的探空资料，解析层结数据，判断是否有效观测，计算稳定度参数，产生图形输出。

第三部分是历史气候分析模块。该模块主要是对历史站点数据进行解析、清洗、聚合，完成年度极值统计、长期趋势分析、EW指数计算、代表事件筛选等工作。

第四，机器学习误差校正模块。本模块完成训练样本准备、观测真值对齐、特征提取、模型训练、模型加载和运行时校正的应用，是提高预报可信度的关键部分。

第五章农业气象服务模块。该模块在天气及历史分析结果的基础上，采用作物知识库、积温计算模型等手段，完成农业风险预警、农事建议的输出。

除上述五个核心业务模块外，系统还专门设置了首页导航页作为平台级入口页面。该页面并不承担复杂分析计算，但负责整合实时天气概览、统一导航、智能提醒和各主要模块跳转关系，是用户理解系统整体功能结构和进入后续业务流程的起点。因此，在系统设计中，首页不是可有可无的装饰性页面，而是连接各模块的重要枢纽。

\subsection{4.3 数据流设计}

系统数据流可以简要归纳为外部数据获取、本地清洗和组织、分析计算、接口封装、页面展示这五个过程。

在多模式预报流程当中，系统会先去开放预报接口申请 ECMWF，GFS，ICON 等模式的数据，之后依照统一的目标站点坐标和字段组织准则来创建标准化的输出。对于ECMWF数据，如果模型已经加载并且可以识别出预报起报时次，则系统会调用机器学习校正逻辑，得到校正后的结果，最后将原始结果和校正结果一起返回给前端。

在探空分析流程中，系统先根据页面输入的北京时间换算为标准 UTC 探空时次，再向怀俄明大学接口请求数据。若指定时次无有效层结数据，则允许回退到前一个标准时次；若解析成功，则进一步提取层结表、计算稳定度指数并生成 Skew-T 或其他图形。

在历史气候流程中，系统从本地历史观测文件中读取多年样本，完成日期解析、变量清洗和统计聚合，再生成年度统计、趋势结果和 EW 指数代表事件结构，以供年度页和趋势页展示。

在农业气象流程中，系统则将预报结果与作物知识库结合，计算积温进度、成熟时间估计及风险提示信息，并将结果展示在农业页面中。

\subsection{4.4 站点与时次统一策略}

由于系统同时涉及预报数据、观测数据、历史数据和探空数据，若站点配置和时间口径不统一，容易造成结果解释混乱。因此，系统在设计中强调统一的站点和时次策略。

在站点方面，系统将杭州国家站作为主要目标站点，并在预报服务与机器学习误差校正链路中尽量统一到相同的经纬度配置，以避免出现“训练一个点、运行另一个点”的情况。

时间上将本地时间、UTC 标准时次、预报起报时间、目标预报时间加以区别。探空分析页面输入以北京时间为准，但是后端会转换成UTC并统一格式为00Z或者12Z；机器学习误差校正部分使用issue time和target time进行明确的区分，防止训练集和推理阶段口径不同。

该种统一策略一方面可以提高系统的实现一致性，另一方面也可以给论文的结果解释提供更加可靠的基础。


\section{第 5 章 关键功能模块设计与实现}

在对各个核心业务模块进行展开之前，首先要说明首页导航页在系统中起到的作用。不同于单纯的静态欢迎页，本系统首页具有“统一入口、实时概览、模块分发”这三种功能。一方面，首页用统一的导航栏和功能卡片把多模式预报、探空分析、历史气候、农业气象等模块组织到同一个视觉框架里，使整个平台的结构更清晰；另一方面，首页顶部实时天气区和智能提醒区又能提前给用户在进入专业分析页面之前提供当前天气状态、体感信息以及重点提醒，从而提高平台整体的连贯性和可用性。因此首页并不是全文的独立技术主线，但是它是系统页面组织和业务引导中不可缺少的一环。

\subsection{5.1 多模式预报模块设计与实现}

\begin{figure}[htbp]
\centering
\caption{图 5-1  多模式预报模块流程图}
\label{fig:5-1}
\end{figure}

\begin{figure}[htbp]
\centering
\caption{图 5-2-1  多模式预报页面示意图}
\label{fig:5-2}
\end{figure}

\begin{figure}[htbp]
\centering
\caption{图 5-2-2  多模式预报页面示意图}
\label{fig:5-2}
\end{figure}

多模式预报模块属于本文系统的重要模块之一，也是本文业务主线中最为重要的部分。不同于传统的天气页面只给出一个预报结果，本模块更加强调两个问题，一是怎样使未来天气的变化更加连续、直观，二是怎样让用户看到不同的模式之间的差别，从而理解预报结果的不确定性以及参考价值。

按照此思路，系统并没有把预报功能设计成单一的列表，而是分成了两个互相配合的层次。第一部分为72小时精细化预报，主要给出未来三天温度、降水、湿度、风速在各时间段的变化情况，适合普通用户做短期出行和生活安排使用。第二部分为多模式对比分析，主要用以展示 ECMWF、GFS、ICON 等模式在同一天线上的差异，从而让使用者看到各个模式对于天气过程强弱、转折时间及趋势发展是否一致。用这样的双层结构来保留日常使用方便的同时又提高页面分析能力。

技术实现上，本模块使用Flask作为页面和接口的载体，用Open-Meteo API作为主要的预报数据来源，在服务层对多模式结果进行统一组织。前端使用Plotly.js画图，时间序列温度、降水、风速等可以用交互式的曲线来表现。与静态表格相比，此种方式更能体现连续的天气过程，也更有利于在同一个坐标系中比较不同模式之间的差异。

为了加快页面响应速度，在该模块中使用了短时缓存以及异步加载的方式。页面在第一次加载的时候先完成整体结构的渲染，然后分别请求精细化预报和多模式数据，从而避免由于一个接口响应较慢而阻塞整个页面。同时短周期缓存可以降低对于同一个预报数据重复请求的情况，在保证结果时效性的基础上提高用户的体验。对多模式场景系统采用并行获取不同模式结果的方式来减小等待时间，使页面更适合作为在线访问。

本模块除了原始预报的显示之外，还有对机器学习误差校正的主要应用环境。系统把 ECMWF 预报链路同误差校正模型联系起来，用户在查看逐小时预报的时候可以决定是否采用校正结果。该种设计不是把机器学习模块作为一个独立的页面来实现，而是在预报功能的基础上，将机器学习模块作为一个增强功能嵌入到主业务流程中。因此用户既可以用原始预报为基准，也可以在需要的时候查看订正后的温度、湿度、风速和降水结果，使用体验更自然。

需要指出的是，系统并没有把机器学习校正无差别地应用到所有的模式上，只在 ECMWF 这一条主链路上启用。该处理符合训练口径和运行口径一致的原则，也可以防止把同样的校正结果机械地套用到不同的模式上所造成的误导。同时系统还给校正结果设计了回退机制，当模型不可用或者运行时的条件不符合要求的时候，页面仍然可以稳定地返回原来的预报结果，保证主功能不受影响。

总体上来说，多模式预报模块把 Flask 页面组织、Open-Meteo 数据调用、Plotly.js 图表展示、短时缓存、异步加载、机器学习误差校正有机地结合在一起，不但提高了天气信息的可读性，而且使系统由单纯的预报展示平台发展成为具有一定分析性质的综合气象服务平台。

\subsection{5.2 探空分析模块设计与实现}

\begin{figure}[htbp]
\centering
\caption{图 5-3  探空分析模块流程图}
\label{fig:5-3}
\end{figure}

\begin{figure}[htbp]
\centering
\caption{图 5-4  探空分析页面示意图}
\label{fig:5-4}
\end{figure}

探空分析模块是本系统中专业性最强的功能模块之一。相比地面天气预报主要回答“未来天气怎么样”，探空分析更关注“大气当前处于什么状态、是否具备不稳定条件、是否存在对流发展潜势”。因此，这一模块的作用并不只是补充一个新的数据页面，而是承担了平台专业分析能力展示的重要任务。多模式预报模块解决的是“未来会发生什么”，而探空分析模块解决的是“当前大气为什么会这样，以及是否具备进一步发展的环境条件”。两者结合后，系统不仅能提供面向公众的天气结果，也能提供一定程度的气象诊断依据。

从业务需求出发，系统引入探空模块主要是为了解决三个问题。第一，单纯查看地面温度、降水和风速，并不能完整判断对流潜势、逆温层结构或高空风场特征，特别是在雷雨、短时强降水和强对流天气背景下，仅靠地面信息往往不足以支撑分析。第二，原始探空资料通常以文本和层结表形式提供，专业性较强，普通用户难以直接阅读，因此需要将其转换为图形和关键参数。第三，系统不仅要为专业用户保留探空图和指数，也要为普通用户提供风险等级、通俗解读以及面向实际场景的影响提示。

基于这一思路，探空分析模块被设计成一个集“探空资料获取、垂直结构分析、专业图形展示、关键指数提取和结果解释”于一体的综合诊断模块，而不是单纯的数据下载功能。当前系统的探空数据主要来自怀俄明大学在线探空资料服务，后端通过 requests 请求原始资料，并使用 Pandas 对层结信息进行整理，再交给分析和绘图模块处理。图形部分主要依赖 Matplotlib 和 MetPy 完成专业气象图生成，页面端则采用 Bootstrap 布局和前端 JavaScript 异步更新方式完成展示。

从页面结构上讲，探空分析模块是用左侧控制、右侧展示的方式来安排的。左侧为时间选择、图形类型切换、关键指数显示，右侧为探空图、风险评估、分析说明。该结构使用户可以按照“先选时次、再看图形、再看结论”的顺序来完成阅读，更符合探空资料使用的习惯。目前模块主要有四种图形，T-lnP图、Skew-T图、高空风矢图、卡通直观图。T-lnP图、Skew-T图主要用来表示温度、露点、风场、大气稳定结构；高空风矢图反映不同高度的风速、风向的变化；卡通直观图把专业结果更易理解地转为可视化的形式，方便非专业用户快速掌握大气状况。

除了图形以外，系统还会同步显示一组重要的指数，即CAPE、CIN、K指数、抬升指数、0-6km垂直风切变和整层可降水量。这些指标被集中到参数区，配有简要说明，方便用户在没有完全依靠探空图的情况下先做初步判断。CAPE反映对流有效位能，CIN反映抑制能量，风切变与风暴组织化程度有关。对于非专业用户来说，图形和参数的结合使用比单独看探空图更容易理解；对有基础的用户来说，这些指标可以很快找到分析的重点。

探空模块的另一个特点就是除了给出参数、图形外，还会给出风险等级以及场景化的解释。系统目前会将探空结果进行归纳，给出较为通俗的解释，尝试从航空、农业、日常生活等角度来补充影响提示。该种设计使探空模块不再只是专业图件的展示页面，而是可以将探空分析的结果转化为实际应用信息的综合模块。对航空来说，探空资料可以用来判断层结稳定性、风切变以及对流风险；对农业来说，可以理解强对流、短时强降水或者逆温状况给农业带来的影响；对公众用户来说，可以转化为比较直观的天气风险提示。

探空模块最值得强调的支撑设计之一，就是对时间口径进行统一处理。由于页面面向国内用户，时间输入以北京时间理解更自然；但是探空业务资料一般用00Z、12Z标准UTC时次组织，所以系统在后端会先把北京时间转换成UTC，再规整到最近的标准探空时次。若目标时次没有有效资料，系统可以回退到前一个标准时次，将最终实际分析所用的时次反馈给用户。该机制本身不是模块的主要卖点，但是它可以保证探空分析在实际使用过程中是稳定的、可解释的。

另外，系统对于参数处理更看重结果的可靠性而不是形式的完整。比如在CAPE进行补充计算的时候出现负值的情况，页面不会给出一个可能会误导用户的数值，而是一直用N/A来表示异常。与把所有的指标都用数字来表示的思想不同的是，这样一种处理方式更符合面向真实用户的应用系统的要求，并且有助于减少错误结论的流传。

探空分析模块把怀俄明大学探空资料、requests数据抓取、Pandas结构化处理、Matplotlib和MetPy专业绘图、Bootstrap+JavaScript页面交互结合起来，使系统在多模式预报的基础上又有了垂直大气结构诊断和场景化解读的能力。该模块明显加强了整个平台的专业程度，也是论文中能够体现系统“分析能力”而不是“展示能力”的部分之一。

\subsection{5.3 历史气候分析模块设计与实现}

\begin{figure}[htbp]
\centering
\caption{图 5-5  历史气候代表事件识别流程图}
\label{fig:5-5}
\end{figure}

\begin{figure}[htbp]
\centering
\caption{图 5-6  历史气候年度页示意图}
\label{fig:5-6}
\end{figure}

历史气候分析模块是本系统中最具统计分析特征的功能模块之一。与多模式预报关注未来几天、探空分析关注当前时次的大气结构不同，历史气候模块更强调从多年观测样本中提取背景信息，帮助用户理解某一年的整体气候特征、长期变化趋势以及典型异常事件在历史中的相对位置。正因为有了这一模块，系统才能把“短期天气”“当前环境”和“长期背景”连接起来，形成更完整的分析链条。

从业务需求出发，用户在看到某一年的高温、暴雨或低温事件时，通常不仅关心它本身的数值大小，还会进一步关心这一年整体上是偏暖还是偏冷、偏湿还是偏干，以及某个异常事件到底只是绝对值较大，还是在历史背景中也具有代表性。基于这一需求，系统将历史气候分析组织为三条主线：年度气候概览、长期趋势分析和代表事件识别。这样的设计使历史气候模块不再只是“看历史数据”，而是能够从年度统计、多年趋势和典型事件三个角度理解气候特征。

在页面结构上，历史气候模块当前分为两个主要页面，即年度分析页 /history/yearly 和趋势分析页 /history/trend。年度分析页主要反映某一年的平均气温、总降水量、雨雪日数、高温日等指标，与标准气候态进行比较；趋势分析页将多年年度统计结果连贯起来，用来观察气温、降水、极端天气频次在更长时间尺度上的变化规律。双页面结构既可以满足用户对某一一年的细看需求，也可以满足用户对多年总体变化的把握。

从技术实现角度来讲，历史气候模块用 Flask 作为页面入口和接口组织形式，分析层主要是使用 Pandas 和 NumPy 来完成历史观测数据的解析、清洗、按日聚合和统计计算。可视化部分用Plotly、Plotly Express和前端Plotly.js完成年度页和趋势页中交互式图表的展示，比如月度对比图、温度分布图、气温趋势图、降水趋势图和极端天气频次图等。风玫瑰图同样用极坐标图的方式表现出来，使用户可以更加直观地看出年度或者月度风况的分布情况。

历史气候模块的第一个重要功能，是年度气候概览。系统会对指定年份的年平均气温、年降水量、雨日、雪日、高温日、低温日等指标进行统计，并与标准气候态进行对比，帮助用户快速判断该年份是偏暖、偏冷还是偏湿、偏干。当前系统支持两套气候态口径，即 1981-2010 和 1991-2020。允许用户在不同气候态之间切换，可以提高结果的可比性，也使模块更符合气候分析业务的实际表达方式。

第二个功能就是长期趋势分析。系统把多年的年度统计结果串联起来，主要体现气温变化趋势、降水变化特点、高温日、低温日、暴雨日、雷暴日、降雪日等极端天气频次的改变状况。毕业设计论文中这一部分内容尤为重要，它使模块具有了较强的分析性、研究性，而不仅仅是某一年份的静态说明。

第三个也是最具有特色的功能就是事件识别。系统在这一步并没有只保留传统的“原始极值榜单”，而采用的是“原始极值+EW指数”双模式的设计。原始极值模式是以绝对数值的大小来排列的，如年内最高温、最低温和最大日降水量等，直观易懂，所以被当作默认的显示方式。但是只看绝对值的大小不能完全体现一个事件在历史背景中所具有的特殊性，所以系统又引入了EW指数模式。

EW 指数模式的核心思想是：将某一日的高温、低温或降水事件放回历史同月同日样本中比较，观察它在历史背景中所处的位置。对于高温和降水，百分位越高说明越异常；对于低温，则采用反向百分位处理，使越低的温度对应越高的极端程度。为了避免仅凭同日样本就把绝对值不强的事件抬上榜单，系统还加入了同季绝对门槛限制。与此同时，为避免连续高温、连续低温或持续降水过程在榜单中重复刷屏，系统又进一步采用了过程去重思路，将相邻异常日合并为一个过程，并仅保留其中最具代表性的一天。最终，EW 结果按春、夏、秋、冬四季分组展示，而不是全年混排，从而避免某一季节占满全部榜单。

这套“历史同日背景比较+同季门槛+过程去重+季节分组”的设计，把历史气候模块由简单的图表展示提升到代表事件识别的层次上。对用户来说，它能够更好地回答“这一年哪些天气过程最值得被记住”，对于论文来说，它也是一份很好的模块分析亮点。

从总体上讲，历史气候分析模块把历史观测数据统计、气候态对比、长期趋势分析、EW代表事件识别、Plotly系列图表展示有机地结合在一起，具有较强的长期背景分析能力。它既弥补了平台对于时间尺度的不足，又加强了系统对于异常天气以及气候变化问题的解释力度。

\subsection{5.4 机器学习误差校正模块设计与实现}

\begin{figure}[htbp]
\centering
\caption{图 5-7  机器学习误差校正流程图}
\label{fig:5-7}
\end{figure}

机器学习误差校正模块与多模式预报模块的关联最强，它并不是一个孤立存在的算法实验功能，而是直接服务于 ECMWF 预报结果质量提升的增强链路。对用户而言，真正关心的仍然是“未来天气会怎样”，而这一模块所做的事情，是在原始预报结果基础上尽量减少系统性偏差，使页面展示的温度、湿度、风速和降水更贴近真实观测。

从业务背景来看，结合江南地区特别是杭州样本的使用体验和本系统训练、测试样本的表现，可以看到 ECMWF 在中远期时效上往往存在一定的偏暖、偏湿以及降水量级偏差倾向。这种系统性误差并不意味着模式本身不可用，而是说明其原始输出在落到具体站点业务场景时，仍然有必要结合历史样本进行后处理订正。也正因为如此，本模块的重点并不只是“用了机器学习”，而是“如何把模型真正接入预报业务链路，并让结果更可信”。

在数据基础上，当前系统的训练样本主要来自 data/hangzhou\_openmeteo 目录下持续采集的新版 ECMWF 预报文件。按照训练清单目前共对100个预报文件进行了审计，其中97个进入到训练流程当中，另外三个由于采集的时间不符合支持窗口的要求而被剔除出去。训练和运行时都使用杭州国家站/馒头山作为目标站点，坐标为30.2444，120.1528。该种统一的动作十分关键，它很好地解决了旧方案里预报点、训练样本以及观测点的口径不完全相同所造成的难题，使得模型的输出更加具有业务价值。

在样本组织上，系统围绕 issue\_time 进行统一处理。对于新版 HZ\_forecast\_*.csv 文件，凌晨 04:15 左右采集的文件映射到前一日 20:00 本地起报，对应 12Z 周期；下午 16:15 左右采集的文件映射到当日 08:00 本地起报，对应 00Z 周期。训练、验证和测试的划分也以 issue\_time 为主线，而不是对所有目标时刻随机打散，从而避免信息泄漏。与此同时，运行时预报页面也会显式传入 issue\_time，使训练与推理采用的是同一套时次口径。

系统用scikit-learn提供的随机森林模型来建立订正链路。温度、湿度、风速用回归模型做误差订正，降水用两阶段法，先判断有无降水，再估计有降水时的降水量。特征设计的时候系统地保持了预报值本身、时间特征（小时、月份、工作日/周末）、时效特征（lead\_hours、lead\_category、周期性时间编码）这些。由此可以看出本模块并没有一味地追求复杂的模型，而是在目前的业务场景下选择了更合适、也更易于和运行时接口进行对接的方案。

在业务接线上，机器学习误差校正模块并没有单独做成一个独立页面，而是直接嵌入多模式预报模块，作为 ECMWF 预报结果的增强能力出现。当前前端在启用 ML 校正时，会将逐小时 ECMWF 预报数据与运行时估计的 issue\_time 一并提交给 /apply-ml-correction 接口；后端完成校正后，再将温度、湿度、风速和降水的订正结果返回页面。如果模型已成功加载，则页面显示校正结果；如果模型未加载、时次不明确或特征条件不满足，则系统自动回退到原始预报值。这样一来，模型不会打断预报主链路，而是以“可开关、可回退”的方式增强 5.1 模块。

当前系统不仅完成了模型训练，还同步生成了 training\_manifest.json 和 metrics\_report.json 两类追踪产物。前者记录训练数据目录、站点信息、文件审计情况、样本时间范围，后者记录连续变量和降水分类/回归的交叉验证和最终测试结果。这就使得模型训练不再只是单次的离线实验，而成了可以追溯、可以复现、可以不断迭代的一个业务过程。

为了更好的体现两版模型在特征设计方面存在的不同，本文将V1和V2在特征设计上主要的不同点归纳成表5-1。

V2是在V1的基础上，加入了季节、起报周期、短期变化量这些特征以更好的表现误差的时序依赖以及季节性规律。相比于只依靠基本预报值和通用时间变量，这类增强特征更加重视同一模式误差随着季节、起报时次和预报时效一起变化的业务事实，因此更适合用作后面boosting实验方案的输入基础。

表 5-1  V1 与 V2 机器学习误差校正特征差异对比

从最终测试结果看，温度、湿度和风速都表现出不同程度的改善，其中温度和风速的提升更为明显。下表给出连续变量在最终测试集上的主要评估结果。

表 5-2  V1 误差改进幅度

对于降水，系统同时给出分类结果与量级结果。需要诚实说明的是，降水校正并非在所有分类指标上都优于原始预报，但在降水量级误差上改善较为明显，这也说明降水后处理的难度高于温度和风速，不能只依赖单一指标下结论。下表给出降水在最终测试集上的主要评估结果。

表 5-3  V1 降水误差与原始版对比

这一结果是可以在论文中作出合理解释的。降水本身就具有更强的随机性、局地性，即使是在同一个模式、同一个站点口径下，降水是否出现也会比温度、风速更难被准确刻画。其次，目前校正模型策略上更多地减少了“报了雨但实际上没有下雨”的误报，因此Recall有所降低，Precision基本不变或者略有提高，说明模型整体上更保守。从业务应用角度来讲，降水量级误差明显降低有实际意义，当系统判定有降水的时候，校正后的雨量值一般更接近真实的状况，对农业灌溉安排、水库泄水调度、田间排涝准备等场景来说，比起只注重提高降水召回率而言，更有实际的应用价值。

需要指出的是，本模块的意义并不在于保证所有变量、所有折次都一定优于原始预报。事实上，交叉验证中的不同折次仍可能出现改进幅度波动，甚至个别折次变差的情况。但在更统一的站点口径、更统一的时次口径以及更严格的评估条件下，这一版方案显然比旧版更真实，也更具参考价值。对于论文而言，这种“结果更可信”本身就是非常重要的结论。

因此，机器学习误差校正模块把scikit-learn、Pandas、NumPy、Flask、训练清单和评估报告有机地结合在一起，使得系统在多模式预报的基础上又有了后处理优化的能力。它不但提高了主要预报结果的实用性，而且使整个平台工程完整性、技术深度都得到了提高。

\subsection{5.5 农业气象服务模块设计与实现}

\begin{figure}[htbp]
\centering
\caption{图 5-8  农业气象服务模块流程图}
\label{fig:5-8}
\end{figure}

\begin{figure}[htbp]
\centering
\caption{图 5-9  农业气象总览页示意图}
\label{fig:5-9}
\end{figure}

农业气象服务模块是本系统面向行业场景延伸的重要组成部分，其核心目标不是简单罗列若干作物信息，而是把前文已经构建的天气预报、指标计算和风险识别能力，进一步转化为面向农业生产的管理建议。相较于公众天气服务更关注“未来天气如何变化”，农业场景更关心“这种天气会对当前作物阶段产生什么影响，以及接下来该做什么”。因此，本模块围绕“作物状态识别、环境适宜度判断、风险预警与农事建议”四条主线组织业务逻辑，使系统能够从通用气象平台进一步延伸为具有一定决策支持能力的农业气象应用。

在业务对象上系统选择了浙江地区具有代表性的一些作物，即单季晚稻、西湖龙井、杨梅、柑橘和小白菜，并且给每种作物建立了包含生育阶段、适宜温度、水分需求、积温参数以及气象风险阈值的作物知识库。不同作物的管理重点也不一样，茶叶更注重倒春寒、高温灼伤，杨梅更注重成熟期暴雨、大风、高湿引起的落果烂果风险，柑橘更注重冻害、夏季日灼，小白菜则具有周年可种、周期短等特点。将业务知识结构化后，在天气数据出现的时候，系统不仅可以给出气温多少、降水多少等天气情况的描述，还可以进一步解释这些天气条件给不同的作物带来了怎样的影响。

在页面结构上，农业模块由总览页 /agro-dashboard 和作物专题页 /crop/<crop\_id> 组成。总览页是全局农业气象研判的入口，集中显示风险提示滚动条、作物实时状态列表、未来七天农事日历、风险预警清单和全局农情研判结果，使用户在开始使用时就对近期农业生产环境有一个大致的认识。作物专题页进一步细分到单个作物上，呈现当前生育阶段、积温进程、阶段性风险阈值、未来 7 天环境变动以及个性化的农事提议，让用户可以针对某个作物展开更为精确的管理判断。该种“总览+专题”的双层结构，既可以满足管理者快速浏览整个情况的需求，也可以满足种植者对单个作物具体建议的详细查看。

从业务处理流程看，农业模块以上游预报模块输出的未来天气数据为基础，由后端先将逐小时预报转换为逐日农业气象信息，再进入作物分析流程。系统在 build\_agro\_forecast\_list() 中按日聚合温度、降水、湿度和风速，并进一步计算 ET0、净有效降水 precip\_eff、浅层土壤湿度和短波辐射等指标。这样做的意义在于，农业场景并不总是直接关心小时级的原始天气值，而更关注“未来几天总体偏湿还是偏干、适不适合作业、是否存在持续性不利条件”。通过把天气数据转换成更贴近农业语言的日尺度指标，系统后续的风险判断和建议输出也就更容易与实际管理需求对接。

系统根据日期来判断各个作物目前所处的生育阶段，再计算出各个作物的适宜度评分以及风险预警。适宜度评分由温度、水分、极端天气三者共同决定，得到的综合分值再用在总览页上直观体现各种作物目前所处的环境状况。风险预警是由 agro\_alert\_engine.py 来驱动的，它把未来七天每天的预报和作物阈值进行对比，找出低温冻害、高温热害、暴雨渍涝、大风落果、高湿病害等风险，并且也发现了适合喷药、施肥、修剪等田间作业的时间窗口。因此模块输出的不单单是被动的风险提醒，还有“可利用的天气机会”，更加符合农业管理“趋利避害并重”的实际需要。

积温分析属于本模块的重要功能。根据不同的作物设置基点温度和目标积温，在作物专题页中给出当前积温、未来7天的预测新增积温和总体进度百分比，供作物发育进程及成熟节律参考使用。历史积温部分采用ERA5-Land再分析日值数据，计算得到的当前实现是根据指定起始日期到昨日的日最高、日最低气温，然后转换成累计积温；未来 7 天新增积温直接使用系统当前主链路中的 ECMWF 预报结果进行估算。因此选择这样的组合，一方面是因为农业场景下连续、长期、稳定的单站历史逐日观测很难完整获取，再分析资料在时间连续性以及历史覆盖范围上更符合支撑整季积温累计的要求；另一方面，积温属于长期积分量，日尺度误差在跨月甚至跨季累积时会被一定程度平滑，所以ERA5-Land用来做历史积温估算是合理的。相比之下，未来7天中部分更加重视对当前业务预报的响应，所以仍然使用ECMWF作为前向估计的依据，也可以保持和平台预报主链路一致。对小白菜等短周期作物来说，页面可以按照用户的播种日期来计算出合理的估计。积温功能的价值就在于它把天气条件和作物生长过程联系得更加紧密，使用户看到的不只是单次天气的影响，而是与作物发育进度直接相关的连续过程判断。

就建议生成而言，系统采取“AI 建议 + 规则回退”并行的办法来改善可用性。总览页的全域农情研判是由farming\_ai\_adviser.py结合预警信息、未来天气、作物状态生成的，作物专题页是通过/api/agro/advice接口为单个作物请求个性化的农事建议。如果DeepSeek API可用，系统就会给出更加自然、更加完整的农事说明；如果接口不可用、返回异常或者解析失败，那么就直接采用规则引擎给出的结构化建议。该种设计可以防止模块对单个AI服务产生硬依赖，在真实的运行环境中也可以稳定地给出可执行的提示信息。前端页面使用 Flask 模板、Bootstrap 组件、JavaScript 交互和 Chart.js 图表来实现总览展示、AI 刷新、环境曲线可视化等功能，使农业模块具有较好的业务完整性以及页面表达能力。

从应用价值来说，农业气象服务模块不是全文技术主线，但是也是系统场景落地能力最强的部分之一。它表明本系统并不只是简单的天气展示平台，而可以将气象要素同作物生长、田间管理、灾害防范联系起来，产生更加贴近实际生产的服务输出。对于茶园管理、果园病害防治、水稻抽穗扬花期防灾、蔬菜短周期种植安排等场景，以天气分析为基础、以作物管理为落脚点的服务方式有较强的实用性。该模块既可以体现系统应用的拓展性，也可以体现预报分析能力向行业服务转化的完整闭环。


\section{第 6 章 系统测试与结果分析}

\subsection{6.1 测试目标与方法}

本章测试工作的重点，不只是说明系统页面能够打开或接口能够返回结果，更重要的是验证前文提出的多条业务链路是否真正形成闭环，并能够输出具有参考意义的分析结果。结合当前项目实际情况，系统测试主要围绕四个方面展开：一是主要页面与接口能否稳定访问；二是关键分析逻辑是否符合设计预期；三是机器学习误差校正结果是否真实、可解释；四是在外部数据波动或参数异常时，系统是否具备合理回退机制。

系统测试方式为自动化测试加页面级人工检查。自动化部分主要使用 pytest 与 Flask test\_client 对历史气候、探空分析、预报时次识别等关键逻辑进行验证；页面层面则对 /advanced-forecast、/upperair、/history/yearly、/agro-dashboard 等主要模块进行了功能检查和结果核对。对于机器学习模块，还保留了 training\_manifest.json 和 metrics\_report.json 两类评估产物，用于支撑模型训练过程与测试结果的追溯。

\subsection{6.2 主要功能流程验证}

从整体流程上看，系统首页可以作为统一的导航入口，把用户引导到多模式预报、探空分析、历史气候、农业气象等主要模块上。各个模块虽然业务侧重点不同，但是都以同一个杭州站点口径、统一的页面导航结构以及相对稳定的的时间组织规则为基础，从而使得整个系统具有较好的整体性。

多模式预报模块测试的重点是72小时精细化预报、多模式对比和ECMWF机器学习校正链路能否协同工作；历史气候模块重点检验原始极值和EW指数模式切换、四季代表事件展示、页面渲染稳定度；探空模块重点检验输入时间转换、数据抓取、参数解析、异常值表达的一致性；农业模块重点检验总览页、作物专题页、AI建议接口之间能否形成作物状态、风险提示、农事建议的闭环。整体测试结果说明各个主要功能模块都可以按照设计意图实现协同工作，系统主业务流程是连续的。

\subsection{6.3 关键逻辑验证结果}

历史气候模块的验证主要看EW指数代表事件筛选逻辑。自动化测试结果表明，系统可以正确地从历史同月同日的样本中提取出来，对低温事件采用反向百分位处理，对于同一个低温、高温或者降水过程中的连续异常日进行合并，防止同一个过程出现多次上榜的情况。同时EW的结果可以按照春、夏、秋、冬四季分组显示，并且控制每个季节每个要素的代表性事件的数量，说明本模块有一定的统计解释力，而不是只停留在图表的展示上。

探空模块主要的验证就是时次转换以及异常参数的处理。测试结果表明，页面输入的北京时间能够正确换算为 UTC 标准探空时次；当用户输入 2026-04-10 08:00 时，系统可正确解析为 2026-04-10 00:00 UTC。接口元数据中可以返回请求本地时次、请求UTC时次、实际解析时次和是否使用回退。当CAPE出现负值或者计算出错的时候，系统会统一用N/A来处理，不会直接显示负数，从而降低了给用户造成误导的风险。

从预报时次统一下手看，经过相关测试发现，新版HZ\_forecast\_*.csv文件的时次识别规则可以保持稳定工作，凌晨04:15左右采集的样本会被识别成前一天20:00本地起报，对应的周期是12Z；下午16:15左右采集的样本会被识别成当天08:00本地起报，对应的周期是00Z；明显偏离支持窗口的异常样本会被剔除。这就意味着训练样本的组织方式和页面运行时的口径是相同的，没有造成隐性的偏差。

\subsection{6.4 模型评估结果分析}

表 6-1  机器学习误差校正连续变量评估结果

表 6-2  机器学习误差校正降水评估结果

机器学习误差校正模块最终的评价结果说明，更一致的站点口径、更严格的测试划分条件下，连续变量有比较明显的改善。温度MAE由原来的2.355°C降到现在的1.832°C，湿度MAE由原来的9.851\%降到现在的9.085\%，风速MAE由原来的4.413m/s降到现在的2.613m/s。从分时效结果可以看出，短时效内改善的更加明显，后处理模型对短中期预报误差有较好的修正作用。

降水变量的结果就更复杂了。分类指标并没有比原来的预报更好，但是Recall、F1基本不变，Precision基本不变或略有上升，雨样本MAE和全样本MAE都明显降低。由此可以看出目前降水模型策略偏向保守，即减少误报降水，但是对于成功识别出的降水，雨量级估计更加接近实际情况。对农业灌溉、水库调度、排涝准备等场景来说，量级误差的改善仍有较高的应用价值。

为了更直观地比较不同的版本误差校正方案在严格时间隔离测试集上的表现，本文在原始预报、V1 和 V2 三种结果之间又增加了统一对比，见表 6-3。需要说明的是，该表对应最后的测试集使用 issue\_time 留出，测试起报时间为 2026年4月4日到2026年4月10日。

从结果可知，本轮V2实验并没有比V1更好，这是合理的。在目前的数据规模大约为 2 个月的时候，随机森林这种 bagging 模型一般比 boosting 模型要更加稳定，对局部噪声以及样本不平衡性都比较敏感。其次Boosting模型一般需要更多的样本、更多的特征工程以及更加细致的参数调节才能发挥出它的优势，目前的研究阶段主要是集中于站点口径统一、issue\_time规则统一以及严格的时段隔离评估。因此本文的核心结论并不是V2一定比V1好，而是说明在小样本量的情况下，随机森林仍然是误差校正的最佳选择；而和只换模型相比，严格控制时间切分、防止目标时刻泄漏、保证训练和运行时口径一致，往往对结果的可信度影响更大。

表 6-3  不同版本误差校正效果对比（严格时间隔离测试集）

注：加粗为最优结果。V1 在温度、湿度、风速、降水量上均优于 Raw 和 V2。

\subsection{6.5 结果讨论}

综合本章测试可以看出，本系统已经形成较完整的业务闭环。多模式预报、探空分析、历史气候、机器学习校正和农业气象服务并非孤立模块，而是在统一站点口径、统一页面入口和统一时间规则下被组织到同一平台之中。测试结果说明，这种组织方式在当前实现中是可运行且基本稳定的。

另一方面，系统测试也表明，本项目真正难的是页面本身，而不是时间规则、统计逻辑和异常处理等容易被忽略的细节。issue\_time 推断、北京时间到UTC的转换、EW榜单去重、CAPE异常值处理等都会影响结果的可信度。目前这些重要细节都已经得到有针对性的验证，所以本系统的结果比以前更加真实、有参考价值。


\section{第 7 章 总结与展望}

\subsection{7.1 全文工作总结}

本文围绕“基于 Flask 的智能气象信息服务系统设计与实现”这一主题，完成了一套面向杭州地区的综合气象信息服务平台设计与开发。与传统仅提供单一天气展示的页面不同，本系统将多模式预报分析作为主线，把探空诊断、历史气候统计、机器学习误差校正和农业气象服务组织到同一平台之中，形成了从未来预报展示、专业环境分析、历史背景支撑到行业应用延伸的较完整业务闭环。

本文使用Flask构建统一的Web应用框架，用模板页面和接口相结合的方式来组织前后端交互。系统按照路由层、服务层、分析层和模型层对系统模块进行拆分，使得多模式预报、探空分析、历史气候、农业服务等各个功能在统一入口下可以协同工作，又可以保持自己独立的数据处理、分析逻辑。该种实现方式既满足了毕业设计系统完整性要求，又为以后继续增加站点、模式、行业等留下余地。

系统对Open-Meteo接入了ECMWF、GFS、ICON等模式数据，可以做72小时的精细化预报以及多个模式的对比。用户可以查看未来逐小时气象要素的演变，也可以用多模式并列结果来观察各个预报方案之间的一致性以及分歧程度，从而提高对天气过程强弱和转折时段的判断能力。本系统的这个主要模块又引入了 ECMWF 的机器学习误差校正链路，使得预报分析从原来的仅仅是对原始模式结果的展示变成了有一定的后处理优化能力。

系统可以对怀俄明大学探空资料进行层结解析、稳定度参数计算、专业图形展示。通过 Matplotlib 与 MetPy 生成的 Skew-T、T-lnP 等图形，系统能够从垂直方向刻画大气温湿风结构，并通过 CAPE、CIN、K 指数、Lifted Index、可降水量和风切变等参数辅助判断对流潜势和稳定度环境。该模块使系统具备了比常规天气页面更强的诊断能力，也使平台可以为强对流天气分析、农业防灾、航空气象保障等更专业的应用提供支持。

历史气候分析能力系统完成年度页、趋势页、代表性极端事件识别逻辑的实现。本文除了保留原始极值榜之外，又加入了基于历史同月同日经验百分位的EW指数，并且用季节门槛、过程去重和四季分组的方式呈现，使得年度极端事件筛选的结果比较符合统计解释性以及业务可读性。这样历史气候模块就不再是简单的展示几张气候图表，而是可以给当前年份的异常天气过程提供更有力的历史背景支持。

本文选取 ECMWF/Open-Meteo 预报链路作为研究对象，在此基础上对杭州国家站/馒头山站点口径重新组织训练样本，统一了训练期和运行期的站点配置、issue time 识别规则以及评估划分方式。新模型对于温度、湿度、风速等连续变量的误差改善比较明显，在风速订正方面改善幅度较大。虽然降水分类指标并没有在所有的维度上都比原始预报好，但是降水量级的误差明显减少，说明该模块有比较好的业务应用价值。更重要的是，该部分工作使系统预报后的处理过程更加真实、统一、可追溯。

在农业气象服务上，系统把天气预报、历史背景、作物知识库三者融合起来，创建起服务于农业生产场景的模块。系统可以针对水稻、茶树、杨梅、柑橘、小白菜等主要农作物给出积温进度预估、适宜性评定、灾害预警及农事建议的输出。历史积温用ERA5-Land再分析资料估算，未来短期积温增量用ECMWF预报链路推算，使农业模块既能考虑长期背景又能考虑短期决策。由此可知，本系统不但具有综合天气分析的能力，而且也有望向具体行业服务转化。

综合来看，本文完成的主要工作并不在于提出全新的算法，而在于围绕真实业务场景，把多源气象数据、统计分析方法、图形可视化和误差校正能力组织成一套可运行、可展示、可解释的综合系统。对于毕业设计而言，这种“以系统实现为主、以分析能力为核心、以业务闭环为目标”的完成方式具有较强的实践意义。

\subsection{7.2 当前系统的不足}

虽然本文已经完成较为完整的系统实现，但是从实际运行情况来看，目前系统还存在着一些不足，主要表现在数据源稳定性、样本数量、结果泛化能力以及部分模块工程鲁棒性等各方面。

第一对外部的数据源还存在很强的依赖性。不论是Open-Meteo预报接口、怀俄明大学探空服务还是历史观测和再分析资料接口，它们的可用性以及响应稳定性都不是由本地系统所决定的。当外部接口出现访问异常、字段变更或者时次缺少的时候，尽管系统可以依靠缓存、回退以及异常提示来缩减影响，但是依然不可避免地会对结果的完备性以及实时性产生干扰。因此目前的系统在“平台能力自主可控”上还存在不足。

其次，机器学习误差校正模块所用的数据样本量太小。虽然本次用新版杭州站ECMWF预报样本进行统一口径重训，训练和测试的程序也比之前版本要更严格一些，但是整体样本时段还比较短，还不足以覆盖各个季节、各种天气类型以及较长的时长下复杂的误差特征。降水变量本身随机性强、局地性强，在有限样本下很难同时达到分类表现和量级回归的表现，所以目前模型更应该被看作是“初步可用的业务增强工具”，而不是成熟的普适订正方案。

再次，系统目前仍然以杭州单站以及周边区域为场景，泛化能力较差。不论是历史气候分析的口径、机器学习误差校正模型，还是农业气象模块中作物知识的配置，都是以杭州地区为研究对象。聚焦可以提高系统的实现完整性，但是也意味着如果把系统直接迁移到其他地区，很多站点参数、阈值规则、作物结构和训练样本都要重新适应。因此，目前版本更适合用作“区域化综合气象服务系统原型”，不能直接推广到全国任意地区。

另外一些分析逻辑虽然已经满足了目前业务的需求，但是仍然存在一定的简化性。历史气候模块中的EW指数仍然以经验百分位、规则约束为基础，探空模块中部分指数在异常情况下还要依靠回退、空态处理，农业模块中的适宜度、风险阈值也大多依据知识规则、经验来设定。这些做法对于毕业设计系统是合理的，但是要使业务化程度更高，还需要在数据驱动建模上更加细致，在业务验证上更加严格。

系统对于文档化、长期维护的支持性还比较弱。目前项目已经形成了一套比较完整的页面、接口、测试、模型成果，但是对于部署说明、数据来源说明、异常场景处理清单以及后续维护规范，还没有形成一个系统化的整理。对于毕业设计来说，该问题并不会影响主要成果的展示，但是从长远的工程演进来看，文档化不足会增大后续的迭代成本。

\subsection{7.3 后续展望}

针对以上不足，本文认为本系统以后可以从数据、模型、功能、应用四个方面继续前进。

第一，从数据层面加大数据源的稳定性、冗余性。对数值预报、探空、观测数据可以创建比较明确的本地缓存、归档系统，逐渐降低对实时外部接口单点可用性的依靠，对于重要的站点可以增加一些权威的、长时间的数据作为本地观测资料，来提高历史分析以及机器学习训练的基础质量。如果将来条件允许的话，可以逐渐形成统一的数据清单和元数据管理机制，使系统数据来源、时段覆盖、使用边界更清楚。

第二，在模型上可以继续扩大机器学习误差校正的样本量和建模深度。随着越来越多的ECMWF历史预报文件以及更长时间的观测真值被收集起来，之后就可以重新训练更加稳定的订正模型，并且按照季节、按时效或者按照天气类型来进行分组建模，从而更好地适应江南地区中远期预报中经常出现的偏暖、偏湿等系统性误差。除了随机森林之外，还可以研究梯度提升树等更适合表格数据的模型方案，但是前提必须是保持站点口径、时次规则的一致性，不能因为算法复杂度的提高而降低结果的可信度。

第三，系统可以进一步发展成为功能更加专业的、更加完备的分析平台。在多模式预报上可以加入更多的集合信息展示以及模式一致性诊断，在探空上可以补充更多的参数产品、典型天气型识别和图文联动的解释，在历史气候上可以丰富趋势检验、分布特征和极端事件对比的能力，在农业上可以增加更多的作物种类、病虫害气象风险和分阶段农事管理模板。由此系统可以由原来的“综合查询与分析平台”发展成现在的“专题化服务平台”。

第四，从应用角度来讲，系统以后会加强智能分析和自动化服务的功能。目前农业模块已经采用AI建议和规则回退并行的方式，在此基础上，也可以把这种思路应用到多模式预报摘要、探空诊断提示、历史气候年度报告自动生成等场景中，使得系统不仅能给出原始的图表、参数的结果，还能给出更加符合用户决策语言的解释性文字。如果再加上自动化报表和定时任务的能力，那么系统就更有可能成为一个面向农业管理、地方气象科普或者专题业务分析的轻量级服务平台。

本文所完成的系统证明了以多模式预报分析为主、以历史气候和探空诊断为支撑、以机器学习后处理和农业服务为延伸的平台建设思路是可行的。未来工作重点不能只是堆砌新的功能，而应该从数据质量、结果可信度、区域适配性、场景化服务深度四个方面进行持续改进。只有在各个方面不断完善之后，系统才可能由毕业设计原型逐渐发展成一个更加稳定、实用的智能气象信息服务平台。



% --- 参考文献 ---
\section*{参考文献}

\begin{enumerate}
\item Grinberg M. Flask Web Development: Developing Web Applications with Python[M]. 2nd ed. Sebastopol: O'Reilly Media, 2018.
\item Harris C R, Millman K J, van der Walt S J, et al. Array programming with NumPy[J]. Nature, 2020, 585(7825): 357-362.
\item McKinney W. Data Structures for Statistical Computing in Python[C]//Proceedings of the 9th Python in Science Conference. Austin: SciPy, 2010: 56-61.
\item Hunter J D. Matplotlib: A 2D graphics environment[J]. Computing in Science \& Engineering, 2007, 9(3): 90-95.
\item May R M, Goebbert K H, Thielen J E, et al. MetPy: A Meteorological Python Library for Data Analysis and Visualization[J]. Bulletin of the American Meteorological Society, 2022, 103(10): E2273-E2284.
\item Munoz-Sabater J, Dutra E, Agusti-Panareda A, et al. ERA5-Land: A state-of-the-art global reanalysis dataset for land applications[J]. Earth System Science Data, 2021, 13(9): 4349-4383.
\item Zippenfenig P. Open-Meteo.com Weather API[CP/OL]. Zenodo, 2024. https://zenodo.org/records/14582479.
\item University of Wyoming. Upper Air Radiosonde Archive[EB/OL]. https://weather.uwyo.edu/upperair/sounding.shtml.
\item Glahn H R, Lowry D A. The use of model output statistics (MOS) in objective weather forecasting[J]. Journal of Applied Meteorology, 1972, 11(8): 1203-1211.
\item Breiman L. Random forests[J]. Machine Learning, 2001, 45(1): 5-32.
\item Pedregosa F, Varoquaux G, Gramfort A, et al. Scikit-learn: Machine learning in Python[J]. Journal of Machine Learning Research, 2011, 12: 2825-2830.
\item Ben Bouallegue Z, Clare M C A, Magnusson L, et al. The rise of data-driven weather forecasting[J]. Bulletin of the American Meteorological Society, 2024, 105(6): E864-E883.
\item 唐卫亚, 等. 天气学分析基础[M]. 北京: 气象出版社, 2016.
\item 肖金香, 叶清, 吴仁烨. 农业气象学[M]. 第3版. 北京: 高等教育出版社, 2023.
\item 王春乙, 张继权, 霍治国, 等. 农业气象灾害风险评估研究进展与展望[J]. 气象学报, 2015(1): 1-19.
\item 中国气象局. 农业气象灾害风险区划技术导则: QX/T 527-2019[S]. 北京: 气象出版社, 2019.
\item 王春乙. 重大农业气象灾害研究进展[M]. 北京: 气象出版社, 2007.
\item 王春乙, 张雪芬, 赵艳霞. 农业气象灾害影响评估与风险评价[M]. 北京: 气象出版社, 2010.
\item 王春乙, 王石立, 霍治国. 近十余年来我国农业气象灾害研究进展[J]. 气象学报, 2005, 63(5): 659-668.
\item 李俊, 杜钧, 刘羽. 北京"7·21"特大暴雨不同集合预报方案的对比试验[J]. 气象学报, 2015, 73(1): 50-71.
\end{enumerate}

% --- 致谢 ---
\section*{致谢}

本论文及相关系统全部由我独立设计、开发与撰写。

首先感谢大学四年里教过我的老师。正是你们在课程教学、实验指导、日常答疑中所传授给我的专业知识以及工程思维，
才使我在完成本课题的过程中打下了良好的基础。

感谢杭州地区气象开放平台、Open-Meteo、怀俄明大学探空数据库等提供公开的气象数据。
缺少这些高质量、容易获得的真实数据，本系统多模式预报、探空分析、历史气候模块就无法运行。

系统开发、论文写作过程中遇到很多意料之外的技术问题，数值预报数据时次对齐、机器学习校正训练推理口径统一。
最终这些困难都会被一一克服，靠的是不断的查阅官方文档、反复地调试代码、不轻易放弃。

最后要感谢我的家人和朋友，在我遇到瓶颈的时候给予我理解、鼓励和默默的支持，使我能安心地完成学业和本论文。

由于本人水平有限，论文中难免有不足之处，希望各位评阅老师、答辩专家提出宝贵意见。

\end{document}
